\documentclass[11pt]{article}
\usepackage{amsmath,amssymb}
\title{The first original nonlinear coefficient after the signed constant-flux bound}\author{}\date{}
\begin{document}\maketitle
\paragraph{What the completed quadratic estimate does and does not give.}
The mixed boundary magnetic loop bound calculation, together with the constant flux spin and magnetic reduction calculation,
bounds the complete original leading quadratic constant-flux response by
$C\log(2+R)E_R/B$, with all relative boundary and gauge determinants.
Here the next coefficient is isolated in the ORIGINAL compact integral,
rather than replaced by an unspecified Gaussian error.

Fix a finite cube $R$ and the exact constant-flux boundary path with
physical angle $\epsilon$. Let $U^*(\epsilon)$ be its unique nearby
tree-fixed action minimum. Use the exact Haar change
$U_j=U_j^*(\epsilon)\exp(i y_j\cdot\sigma)$ on its seven or more independent
chords, as appropriate to $R$. If $S_0$ is the original action without $B$,
write its genuine local Taylor expansion
\[
 S_0(U^*\exp(iy))=v(\epsilon)+\tfrac12 y^TH_\epsilon y
       +T_{3,\epsilon}(y)+T_{4,\epsilon}(y)+\cdots ,
\]
where $T_k$ includes the factor $1/k!$. With $y=z/\sqrt B$,
\[
 BS_0=Bv+\tfrac12z^TH_\epsilon z+
 B^{-1/2}T_{3,\epsilon}(z)+B^{-1}T_{4,\epsilon}(z)+\cdots .
\]
These tensors are in the actual right minimum frame, so their
$\epsilon$ derivatives retain the minimum response and acceleration.
Normalized $SU(2)$ Haar measure on each chord gives
\[
 j(z/\sqrt B)=1+B^{-1}J_2(z)+O_R(B^{-2}|z|^4),\qquad
 J_2(z)=-\tfrac13\sum_j|z_j|^2 .
\]
It is independent of $\epsilon$ before taking Gaussian expectations.
The tree gauge has normalized volume one; adding a further ghost
determinant to this tree integral would double count the quotient.

\paragraph{Actual fixed-size differentiated expansion.}
Let $\gamma_\epsilon$ have covariance $H_\epsilon^{-1}$ and set
\[
 \mathcal C_R(\epsilon)=
 E_{\gamma_\epsilon}T_{4,\epsilon}
 -\tfrac12E_{\gamma_\epsilon}T_{3,\epsilon}^2
 -E_{\gamma_\epsilon}J_2.                                      \tag{1}
\]
For each fixed $R$, uniformly on a sufficiently small fixed interval of
$\epsilon$, the ORIGINAL compact integral satisfies
\[
 F_R(\epsilon):=-\log Z_R(\epsilon)-Bv(\epsilon)
 =c_{R,B}+\tfrac12\log\det H_\epsilon
       +B^{-1}\mathcal C_R(\epsilon)+O_R(B^{-2}),                  \tag{2}
\]
with the same bound after up to two $\epsilon$ derivatives.
Here $c_{R,B}$ retains every fixed canonical/Haar normalization and is
independent of $\epsilon$. Any fixed-only boundary plaquette action is
present in both $Z$ and $Bv$ and cancels only in this stated difference.

For clarity, (2) is a fixed-dimensional Laplace statement, not a
measure replacement: compactness and the unique flat minimum give a
uniform positive action gap outside a small neighborhood for small
$\epsilon$. Inside it the positive Hessian is uniform. A symmetric
cutoff, Taylor expansion through sufficiently many orders, and Gaussian
domination justify termwise integration and two parameter derivatives.
The complement is exponentially small, with derivatives costing only
fixed powers of $B$. Odd polynomial coefficients vanish in the centered
Gaussian. The order $B^{-1}$ term of the original integral is
$E(J_2-T_4+T_3^2/2)$; taking minus its logarithm proves (1).
Expanding two further orders gives the stated differentiated remainder.
Every constant and neighborhood in this justification may depend on $R$.

For canonical $\epsilon=\zeta/\sqrt B$, (2) therefore gives the actual
fixed-block response at zero:
\[
 \boxed{\quad
 F_{R,\zeta\zeta}(0)
 =B^{-1}\left(\tfrac12\log\det H_\epsilon\right)''_{\epsilon=0}
   +B^{-2}\mathcal C_R''(0)+O_R(B^{-3}).\quad}                     \tag{3}
\]
Thus the first POSSIBLE non-Gaussian coupling correction is $B^{-2}$,
not another independent leading $B^{-1}$ term, for this fixed-size
family. This is not a uniform growing-block bound.

\paragraph{The first explicit nonlinear current on the complete $R=2$ cube.}
Every plaquette in this cube contains at most two free links: all eight
free links meet the unique interior vertex. At the flat record the
scalar part of a product of at most two fluctuation exponentials is
even under simultaneous reversal of the fluctuations.
Consequently $T_{3,0}\equiv0$, even before the temporal tree edge is
fixed. The derivative $T'_{3,0}$ need not vanish: the varying
nontrivial external/minimum factors destroy that two-factor reversal
at nonzero $\epsilon$.

Differentiating (1) at this flat record gives the exact finite Gaussian
coefficient
\[
 \mathcal C_2''(0)
 =\left.\frac{d^2}{d\epsilon^2}
       E_{\gamma_\epsilon}(T_{4,\epsilon}-J_2)\right|_0
       -E_{\gamma_0}(T'_{3,0})^2 .                               \tag{4}
\]
All derivatives of the Gaussian covariance in the first term are
retained. In the second term those derivatives multiply $T_{3,0}$
and hence vanish. Formula (4) displays the first nonlinear source
covariance and its quartic/Haar partners; dropping either side is not
a signed estimate. The exact rational product check
the finite algebraic calculation evaluates
$T'_{3,0}=-3/8$ on one explicitly stored rational chord fluctuation.
Thus this cubic source polynomial is genuinely nonzero and
$E_{\gamma_0}(T'_{3,0})^2>0$. This does not assert that the total signed
combination in (4) is nonzero.

On larger cubes a plaquette may contain three or four free links.
At the flat record its cubic action vertex contains
\[
 \sum_{k<l<m}(s_kz_{e_k}\times s_lz_{e_l})\cdot s_mz_{e_m}.
\]
It is absent in the $R=2$ geometry but generally survives. Then the
full derivative of $E_{\gamma_\epsilon}T_{3,\epsilon}^2$ in (1),
including its covariance derivatives, must be used instead of (4).

\paragraph{Exact remaining mathematical input.}
A growing-block conclusion requires a signed physical-energy estimate on
$\mathcal C_R''(0)$ AND a uniform differentiated remainder bound for the
original integral, or a direct original-law resummation controlling the
equivalent residual (5) of the signed gauge quotient and unsummed current calculation. Fixed-dimensional analyticity proves neither.
The previous $R^{14}/B$ actual estimate remains available only in its
admitted mesoscopic range. No size-independent perturbative expansion,
random-exterior admission, global-mode replacement, or S94 follows from
(2)--(4). The completed leading logarithm is a prerequisite, not the
nonlinear comparison.
\end{document}
